\section{Related Work}
\label{sec:related-work}

\paragraph{Self-supervised representations and their hierarchy.}
Contrastive, self-distillation, and joint-embedding objectives
\citep{caron2021dino,oquab2023dinov2,assran2023ijepa,he2022mae} produce
visual features that transfer well to recognition and, increasingly, to
world modeling \citep{zhou2024dinowm}, but standard probing only asks
whether a concept is linearly recoverable somewhere in the network
\citep{alain2016probes}. We ask a functional version instead: whether a
concept, once pushed through the physics it is supposed to govern,
reproduces a prediction the model already relies on. The tall-but-thin
hierarchy framing follows \citet{filling-hierarchy-blog}; the specific
mechanism of predicting a hierarchy of teacher representations follows
\citet{constructive-ssl-blog}, which we repurpose for a target of a
different kind, physical parameters mapped through a fixed solver
rather than another learned embedding.

\paragraph{Latent world models.}
Predicting forward in a learned latent space
\citep{ha2018worldmodels,hafner2019planet,hafner2020dreamer,
hafner2023dreamerv3} is the standard recipe for model-based control
from images. DINO-WM \citep{zhou2024dinowm} predicts forward in a
frozen, pretrained latent with a block-causal transformer; \pathA is
this recipe unmodified, so \pathB has an accurate, unremarkable
baseline to be compared against. None of these models expose the
physical parameters underlying the dynamics they predict.

\paragraph{Physical identifiability.}
\citet{latent-wm-identifiability} measures the gap directly: accurate
next-state prediction with almost nothing about the generating physics
linearly recoverable. Its negative result, that ratio-type parameters
are not individually identifiable no matter how good the probe, is a
structural fact about the dynamics that \textsc{PokeWorld}
(\Cref{sec:experimental-setup}) reproduces deliberately, and its
probe-then-$R^2$ methodology is the certificate we adopt once a model
exists to certify. The same gap turns up elsewhere: physical
plausibility is linearly decodable from video diffusion states that
were never given a physics objective \citep{esmati2026invisiblehand},
and preventing representational collapse in a JEPA world model does not
by itself preserve physical state \citep{zeng2026phylatent}. Neither
asks whether the decodable signal is used for anything, which is what
the functional-use evaluation in \Cref{sec:setup-protocol} is for.
Closest in spirit is \citet{mao2024piwm}, which aligns latent states
with physical quantities under weak, distribution-based supervision and
a \emph{learned} dynamics model. What differs is what is fixed: no
physical label of any kind here, and a dynamics model that stays frozen
throughout.

\paragraph{Physics-informed and structured dynamics.}
Interaction networks \citep{battaglia2016interaction} and
physics-informed objectives build structure into the dynamics function
or supervise physical quantities directly, typically learning that
function end to end. Our solver is instead hand-derived, frozen, and
owns zero learnable parameters (\Cref{sec:method}), closer to
differentiable simulation used as a fixed decoder
\citep{jatavallabhula2021gradsim} than to physics-informed
regularization of a learned one, though that comparison too has a
difference in what trains it: gradSim backpropagates from pixels
through a differentiable renderer end to end, while our solver is never
optimized and its target is a predictor's own belief. \citet{li2026ldr}
takes the opposite structural choice from ours, integrating a video
world model's latent transition kinematically within a single latent
path, learning only a higher-order residual there; we keep the physics
computation in a second head that can be turned off and falsified
independently of \pathA.

\paragraph{Self-distillation.}
Training against a stop-gradient target traces to classical
distillation \citep{hinton2015distillation} and appears label-free in
\citep{caron2021dino,grill2020byol,assran2023ijepa}, all predicting one
embedding from another. We reuse the mechanism for a target of a
different kind (physical parameters through a fixed solver, not another
learned embedding), and the stop-gradient plays a different role: not
preventing collapse, but enforcing the one-way dependency that keeps
\pathB from pulling \pathA toward itself (\Cref{sec:method}).

\paragraph{Contact-rich manipulation.}
Push-T \citep{chi2023diffusionpolicy} is a standard contact-rich
manipulation benchmark; we model its block with the standard
quasi-static limit-surface approximation and treat the disc-for-T-block
mismatch as measured rather than hidden
(\Cref{sec:experimental-setup}). The cart-pole solver is evaluated
against DeepMind Control Suite rollouts \citep{tassa2018dmcontrol}.
